% Overleaf compiler: pdfLaTeX. One cropped, A4-width panel (no full-page wrapper).
\documentclass[10pt,border=0pt]{standalone}
\usepackage[T1]{fontenc}
\usepackage[mainfont]{Oswald}
\usepackage{amsmath,xcolor}

\definecolor{panel}{HTML}{1E1E1E}
\definecolor{frog}{HTML}{32DD0B}
\newlength{\panelwidth}
\setlength{\panelwidth}{200mm}
\setlength{\fboxsep}{5mm}
\setlength{\parindent}{0pt}
\setlength{\jot}{2pt}
\newsavebox{\portfolioPanel}
\DeclareMathOperator{\atanTwo}{atan2}
\DeclareMathOperator{\unit}{unit}
\newcommand{\vabs}[1]{\lvert #1\rvert}
\newcommand{\vnorm}[1]{\lVert #1\rVert}
\newcommand{\steptitle}[1]{{\fontsize{12}{14}\selectfont\bfseries #1\par}\vspace{1.5mm}}
\newcommand{\key}[2]{\noindent\(#1\)\enspace #2\par}
\newcommand{\rulegreen}{{\color{frog}\rule{\linewidth}{0.65pt}}\par}
\newcommand{\legendfont}{\fontsize{9}{11.4}\selectfont}

\begin{document}%
\sbox{\portfolioPanel}{%
\colorbox{panel}{%
\begin{minipage}{\panelwidth}
\color{white}\raggedright
\setlength{\fboxsep}{2.5pt}% Compact frame around the position equation.
\rulegreen\vspace{0.6mm}
{\fontsize{17}{20}\selectfont\bfseries AprilTags on a Tilting Hive\par}
\vspace{1.1mm}
{\fontsize{10}{13}\selectfont
The Hive's AprilTags move as it tilts, so a fixed tag map gives the wrong robot pose.
Limelight estimates each tag's 3D pose from its four corners. We use that pose to
locate the \textbf{fixed hinge}. With robot \textbf{roll and pitch set to \(0^\circ\)}
on level tiles, the hinge axis gives us heading and the hinge point gives us
position. We combine agreeing tags from the same image.\par}
\vspace{1.6mm}
{\legendfont
\(T,C,R,W\) denote tag, camera, robot and field frames.
\(T_{AB}=(R_{AB},t_{AB})\) maps points from \(B\) to \(A\)
using rotation \(R_{AB}\) and offset \(t_{AB}\). Distances below are in inches.\par}
\vspace{1.5mm}

\begin{minipage}[t]{0.308\linewidth}
\raggedright
\steptitle{1. Read the 3D tag pose}
\(\begin{aligned}[t]
t_{CT}&=k(z,-x,-y)^{\mathsf T}\\
R_{CT}&=S R_z(\gamma) R_y(\beta) R_x(\alpha)\\
T_{TR}&=(T_{RC}T_{CT})^{-1}
\end{aligned}\)\par
\vspace{1.6mm}
\legendfont
\key{x,y,z}{Raw optical coordinates in metres.}
\key{k=1/0.0254}{Metres to inches.}
\key{\alpha,\beta,\gamma}{Measured tag roll, pitch and yaw. \(R_x,R_y,R_z\) are rotations about the named axes.}
\key{S}{Axis change \((x,y,z)\mapsto(z,-x,-y)\), from right/down/forward to forward/left/up.}
\key{T_{RC},T_{TR}}{Camera mount and robot pose in tag coordinates.}
\end{minipage}\hfill
%
\begin{minipage}[t]{0.307\linewidth}
\raggedright
\steptitle{2. Find the fixed hinge}
\(\begin{aligned}[t]
h_R&=t_{RC}+R_{RC}(t_{CT}+R_{CT}d_T)\\
u_R&=R_{RC}R_{CT}a_T\\
\epsilon&=\arcsin\vabs{u_R\cdot g_R}
\end{aligned}\)\par
\vspace{1.6mm}
\legendfont
\key{d_T,a_T}{Hinge point and signed unit axis in tag coordinates, from the tag ID.}
\key{h_R,u_R}{Same point and axis in robot coordinates.}
\key{g_R,\epsilon}{Unit down in robot coordinates and measured hinge-axis slope. Use \(g_R=(0,0,-1)^{\mathsf T}\) for a level robot.}
\key{g_W,e_x}{Field down \((0,0,-1)^{\mathsf T}\) and signed hinge axis \((1,0,0)^{\mathsf T}\).}
\end{minipage}\hfill
%
\begin{minipage}[t]{0.343\linewidth}
\raggedright
\steptitle{3. Recover the robot pose}
\(\begin{aligned}[t]
\hat u_R&=\unit\big(u_R-(u_R\cdot g_R)g_R\big)\\
B(a,g)&=[\,a,\ g\times a,\ g\,]\\
Q&=B(e_x,g_W) B(\hat u_R,g_R)^{\mathsf T}
\end{aligned}\)\par
\vspace{0.7mm}
\(\boxed{p=H-Qh_R}\quad \psi=\atanTwo(Q_{21},Q_{11})\)\par
\vspace{1.6mm}
\legendfont
\key{\hat u_R}{Horizontal unit hinge axis.}
\key{B,Q}{Hinge/down basis and robot-to-field rotation.}
\key{H,p,\psi}{Known field hinge point, robot position and heading. \(p_z\) is robot height.}
On level tiles, \(Q=R_z(\psi)\) and
\(\psi=-\atanTwo(u_{R,y},u_{R,x})\).\par
\end{minipage}\par
\vspace{1.5mm}\rulegreen\vspace{0.8mm}

\steptitle{4. Combine tags from one image}
\begin{minipage}[t]{0.475\linewidth}
\raggedright
\(\begin{aligned}[t]
D(\theta)&=\atanTwo(\sin\theta,\cos\theta)\\
\epsilon_i&\le 10^\circ\\
\vnorm{p_i-p_j}&\le 6\,\mathrm{in},\qquad \vabs{D(\psi_i-\psi_j)}\le 10^\circ
\end{aligned}\)\par
\vspace{1.6mm}
\legendfont
Use unique IDs from one exposure. \(D\) wraps heading differences.
Filter tag \(i\) by axis slope, then compare pairs \(i,j\).
The largest group around a tag must be unique, form a strict majority
and agree pairwise. One usable tag can stand alone.\par
\end{minipage}\hfill
%
\begin{minipage}[t]{0.49\linewidth}
\raggedright
\(\begin{aligned}[t]
\bar p&=\tfrac1n\sum\nolimits_i p_i,\qquad
\bar\psi=\atanTwo\!\big(\sum\nolimits_i\sin\psi_i,\sum\nolimits_i\cos\psi_i\big)\\
s_p&=\sqrt{\tfrac1n\sum\nolimits_i\vnorm{p_i-\bar p}^{\,2}},\qquad
s_\psi=\sqrt{\tfrac1n\sum\nolimits_i D(\psi_i-\bar\psi)^2}
\end{aligned}\)\par
\vspace{1.6mm}
\legendfont
\(n\) accepted tags contribute equally to position \(\bar p\) and
circular heading \(\bar\psi\). \(s_p,s_\psi\) report RMS disagreement
in position and heading. Gates shown are defaults. Heading spread is reported in degrees.\par
\end{minipage}\par
\vspace{2mm}
{\legendfont\textbf{\textcolor{frog}{Simulation}}
\enspace Four-tag fusion reduced RMS position error from \(1.01\) to \(0.51\) in
under independent translation noise.\par}
\end{minipage}%
}%
}%
\typeout{Hive panel height: \the\dimexpr\ht\portfolioPanel+\dp\portfolioPanel\relax}%
\ifdim\dimexpr\ht\portfolioPanel+\dp\portfolioPanel\relax>148.5mm
  \PackageError{hive-panel}{Panel taller than half an A4 page}{Reduce the content or spacing.}
\fi%
\usebox{\portfolioPanel}%
\end{document}
